% 本文件由 LaTeX 排版 Agent 根据有效 translation.json 自动生成。
% author=Ignatius of Antioch; work_id=0107; title=依纳爵致罗马人书

\chapter{\WorkTitle{依纳爵致罗马人书}}

% structure: p0001 source=h1 target=omitted reason=page-title-already-rendered-by-parent

% structure: p0002 source=h2 target=section reason=relative-heading-rank; base=h2

\section{\Emph{致候}}

\Person{依纳爵}，亦号\Person{塞奥弗鲁斯}，致那赖至高圣父的威能及\Person{耶稣基督}——祂的独生子——获得怜悯的教会；这教会因那按我们\Person{天主}\Person{耶稣基督}之爱而愿万有者所意愿，而蒙受宠爱与光照；这教会也主持罗马地区，堪当\Person{天主}的、堪当尊荣的、堪当极乐的、堪当赞颂的、堪当成就其一切愿望的、堪当被视为圣的，并主持爱德，由\Person{基督}及圣父而得名；我亦以圣父之子\Person{耶稣基督}之名向她致候：致那些按肉体和圣神，都结合于祂每一条诫命的人；他们无可分割地充满\Person{天主}的恩宠，并从一切异己的污秽中被净化；〔我祝愿〕无可指责地、在我们的\Person{天主}\Person{耶稣基督}内，享有圆满的福乐。\par

% structure: p0004 source=h2 target=section reason=relative-heading-rank; base=h2

\section{\Emph{第一章  身为囚犯，我希望见到你们}}

藉著向\Person{天主}的祈祷，我获得了目睹你们那极其堪当尊荣的容颜的特恩，甚至蒙赐了超乎我所求的；因为作为在\Person{基督耶稣}内的囚犯，我希望向你们致候，如果\Person{天主}的旨意认为我堪当达到终局的话。因为开端已安排妥当，只要我能获得恩宠，毫无阻碍地持守我的份直到终局。因为我怕你们的爱德，恐怕它会伤害我。因为你们容易成就你们所愿意的；但你们若宽免我，我便难以达到\Person{天主}那里。\par

% structure: p0006 source=h2 target=section reason=relative-heading-rank; base=h2

\section{\Emph{第二章  不要救我脱离殉道}}

因为我不愿意对你们行事像取悦于人者，而是取悦于\Person{天主}，正如你们也取悦于祂一样。因为我绝不会再有机会达到\Person{天主}那里了；而你们若如今缄默，也再不会配得更好功业的荣誉。因为若你们对我缄默，我便将成为\Person{天主}的；但若你们爱我的肉身，我便要再跑我的赛程。所以，请祈求，不要赐我比成为祭献给\Person{天主}更大的恩惠，趁祭台尚备妥之时；好使你们聚集在爱中，藉著\Person{基督耶稣}向圣父歌唱赞颂，因为\Person{天主}以为我——叙利亚的主教——堪当从东方被遣送到西方。由世界出发到\Person{天主}那里是好的，好使我向祂复活。\par

% structure: p0008 source=h2 target=section reason=relative-heading-rank; base=h2

\section{\Emph{第三章  宁请求我得以殉道}}

你们从未嫉妒任何人；你们教导了别人。如今我希望那些你们在训诲中所吩咐的，〔凭你们的行事〕得以坚固。只求在我这方面赐予内外的力量，使我不仅说话，而且〔真正〕愿意；并且使我不仅被称为基督徒，而且实际被找到是基督徒。因为若我实际被找到是〔基督徒〕，我也可以被称为基督徒，那时才被视为忠信，当我不再显现于世界时。\BibleQuote{因为看得见的是暂时的，看不见的才是永恒的。}因为我们的\Person{天主}\Person{耶稣基督}，如今祂与圣父同在，更被启示在祂的荣耀中。基督教不是仅属沉默的事，也是属〔显明〕的伟大。\par

% structure: p0010 source=h2 target=section reason=relative-heading-rank; base=h2

\section{\Emph{第四章  容我成为野兽的猎物}}

我致信各教会，并向他们全体强调，我将甘愿为\Person{天主}而死，除非你们阻止我。我恳求你们不要向我表示不合时宜的好意。容我成为野兽的食物，藉它们的器皿我将得以达到\Person{天主}那里。我是\Person{天主}的麦子，让我被野兽的牙齿磨碎，好使我被找到为\Person{基督}的纯洁面包。宁可引诱野兽，使它们成为我的坟墓，不留我身体的一点；这样当我〔在死亡中〕安眠后，便不再麻烦任何人。那时我才真正是\Person{基督}的门徒，当世界连我的身体也看不见。为我在\Person{基督}前祈求，使我藉这些工具被找到为〔献给\Person{天主}的〕祭品。我不像\Person{伯多禄}和\Person{保禄}那样向你们颁命。他们是宗徒；我不过是一个被定罪的人：他们是自由的，而我直到如今仍是仆人。但当我受苦时，我便成为\Person{耶稣}的自由人，并在祂内复活得释放。而现在，身为囚犯，我学会不贪求任何世俗或虚浮的东西。\par

% structure: p0012 source=h2 target=section reason=relative-heading-rank; base=h2

\section{\Emph{第五章  我渴望死}}

从叙利亚直到罗马，我与野兽搏斗，无论陆上海上，无论昼夜，被绑在十只豹子身上，我意指一队兵士，他们即使受了恩惠，反而越发凶恶。但他们的恶行更教导我〔行事如\Person{基督}的门徒〕；\BibleQuote{然而我因此也不自以为义。}（\BibleRef{格前4:4}）愿我享用为我预备的野兽；我祈求它们急于扑向我，我也将引诱它们迅速吞噬我，而不像对待某些人那样因为惧怕而不碰它们。但如果它们不愿攻击我，我就强迫它们这样做。请原谅我〔这一点〕：我知道什么于我有益。如今我开始成为门徒。愿没有任何可看见或不可看见的东西嫉妒我，使我不能达到\Person{耶稣基督}那里。让火与十字架，让野兽的成群，撕裂、破碎、骨头的脱节，让肢体的割断，让整个身体的粉碎，以及\Person{魔鬼}的一切可怕酷刑都临到我身上：只要让我达到\Person{耶稣基督}那里。\par

% structure: p0014 source=h2 target=section reason=relative-heading-rank; base=h2

\section{\Emph{第六章  藉死亡我将获得真生命}}

世上的一切快乐，和这地上的一切王国，于我毫无益处。为\Person{耶稣基督}而死，胜过统管地极。\BibleQuote{人纵然赚得了全世界，却赔上了自己的灵魂，为他有什么益处？}我寻求那为我们而死的那一位；我渴望那为我们而复活的那一位。这是为我存留的收益。弟兄们，请原谅我：不要阻止我活着，不要希望我将我留在死亡的状态中；我渴望属于\Person{天主}，不要把我交给世界。容我获得纯粹的亮光：当我到达那里，我必成为\Person{天主}的人。准许我效法我的\Person{天主}的苦难。若有人内心有祂，让他思量我所渴望的，并让他同情我，因为他知道我如何被束缚。\par

% structure: p0016 source=h2 target=section reason=relative-heading-rank; base=h2

\section{\Emph{第七章  渴望死亡的理由}}

此世王子想要带走我，并败坏我对天主的忠诚。所以，你们中任何人（在\Place{罗马}的）都不要帮助他；反而要站在我这一边，即站在天主一边。不要口口声声说耶稣基督，却将你们的渴望寄托于世上的事物。不要让嫉妒在你们中间找到居所；即使我亲临你们当中时劝你们如此，也不可听从我，反而要相信我如今写信给你们的话。因为我虽然活着给你们写信，但我切望死去。我的爱已被钉在十字架上，在我里面再也没有欲火需要满足；却有活水在我内说话，对我说：来，到圣父那里去。我不再贪图可朽坏的食物，也不贪图今生的享乐。我渴望天主的食粮，天上的食粮，生命的食粮，就是耶稣基督的肉体——天主子，后来成为达味和亚巴郎的后裔；我也渴望天主的饮料，即他的血，那是不朽的爱和永生。\newline{}\par

% structure: p0018 source=h2 target=section reason=relative-heading-rank; base=h2

\section{第八章　请你们善待我}

我不再愿意按照人的方式生活，如果你们同意，我的愿望就会实现。那么，请你们也愿意，好使你们的愿望得以实现。我在这封短信中恳求你们：请相信我。耶稣基督会向你们启示这些事，使你们知道我说的是真话。他是一张完全无谎的口，圣父借着他真实地说了话。请为我祈祷，使我得偿所愿。我写信给你们，不是按照肉身，而是按照天主的旨意。如果我受苦，你们就善意地对我；但如果我遭到拒绝，你们就恨了我。\newline{}\par

% structure: p0020 source=h2 target=section reason=relative-heading-rank; base=h2

\section{第九章　为\Place{叙利亚}的教会祈祷}

请在你们的祈祷中记念\Place{叙利亚}的教会，它如今有天主为牧者，而非我。惟有耶稣基督会管理它，你们的爱也将会看顾它。至于我，我羞愧于被算作他们中的一员；因为我的确不配，是其中最末的一个，如同未到产期而生的人。见：\BibleRef{格前 15:8} 但我蒙了怜悯，成为什么人物，只要我能到达天主那里。我的精神向你们致意，也向那些因耶稣基督的名接待我、不只是当作过客而接待我的教会之爱致意。因为那些在道路上不靠近我的教会——我是指按肉身来说的——已经一城接一城地先我而去（迎接我）。\newline{}\par

% structure: p0022 source=h2 target=section reason=relative-heading-rank; base=h2

\section{第十章　结语}

如今我从\Place{斯米纳}通过\Place{厄弗所}人写下这些事给你们，他们实在是最有福的。与我同在的，还有我所亲爱的一位，\Person{克洛库斯}，以及许多人。至于那些为天主的荣耀从\Place{叙利亚}先我往\Place{罗马}去的人，我相信你们认识他们；那么，请告诉他们我即将到来。因为他们都配得上天主和你们；你们理应在一切事上款待他们。我在九月朔日前的第九天（即八月二十三日）写下这些事给你们。愿你们在耶稣基督的坚忍中坚持到底。阿们。\newline{}\par

\SourceRecord{Translated by Alexander Roberts and James Donaldson.；Ante-Nicene Fathers；Vol. 1.；Edited by Alexander Roberts, James Donaldson, and A. Cleveland Coxe.；Buffalo, NY: Christian Literature Publishing Co.,；1885.；<http://www.newadvent.org/fathers/0107.htm>.}
